% Image credits for the photographs and AI illustrations of Book 5 (University, Year 3), Arabic edition.
% Machine-readable license records live in images/book5/CREDITS.md; keep this
% file in sync with frontmatter/image-credits-book5.tex (one \item per item there).
% As in the Book 3 and 4 Arabic credits: personal names are transliterated, with the
% original spelling in parentheses (the credit itself); institutions, datasets,
% works and licences stay verbatim inside \babelsublr; only the gloss is translated.
\chapter*{مصادر الصور}

الرسوم والخرائط والمخططات في هذا الكتاب أصلية؛ وما رُسم منها من بيانات
منشورة يذكر مصدره في التعليق المرافق أو في هذه القائمة. أما الصور
الفوتوغرافية فمستعملة بمقتضى رخصها الحرة، مع الشكر لأصحابها؛ وروابط المصادر
الكاملة وروابط الرخص لكل صورة مدرجة في \texttt{images/book5/CREDITS.md} في
مستودع الكتاب.

\bigskip
وإلى جانب الصور الفوتوغرافية والرسوم الأصلية، بعض الأشكال رسوم توضيحية
مولَّدة بالذكاء الاصطناعي، أُنتجت بأداة توليد الصور من OpenAI انطلاقًا من
أوامر كتبها محررو الكتاب، ورُوجعت دقتها العلمية قبل إدراجها. وهي لا تُظهر
إلا مشاهد عامة، ولا تمثّل أبدًا مكانًا حقيقيًا أو شخصًا حقيقيًا أو عيّنة
حقيقية. والأمر الكامل لكل صورة مولَّدة مسجَّل في
\texttt{images/book5/ai/PROMPTS.md} في المستودع.

\bigskip
\noindent الصور الفوتوغرافية والبيانات، فصلًا فصلًا:
\begin{itemize}\small
  \item الفصل~1 (علم الزلازل وتصوير باطن الأرض): \babelsublr{PREM}، دزيونسكي وأندرسون (\babelsublr{A.~M.~Dziewonski and D.~L.~Anderson}، 1981)، كما جُدول في \babelsublr{ObsPy}؛ و\babelsublr{ak135}، كينيت وإنغدال وبولاند (\babelsublr{B.~L.~N.~Kennett, E.~R.~Engdahl and R.~Buland}، 1995).
  \item الفصل~3 (الجاذبية وشكل الأرض): القمران التوأمان \babelsublr{GRACE Follow-On}، تصوّر فنّي، \babelsublr{NASA/JPL-Caltech}، ملك عام (\babelsublr{Wikimedia Commons}). ميزان كتلة الغطاءين الجليديين من \babelsublr{IMBIE Team} (2021)، \babelsublr{UK Polar Data Centre}، \babelsublr{Open Government Licence v3.0}.
  \item الفصل~4 (المغناطيسية الأرضية والمغناطيسية القديمة): المجال المغناطيسي الأرضي المرجعي الدولي، الجيل الرابع عشر، \babelsublr{IAGA Division~V Working Group V-MOD}، ومعاملاته توزّعها \babelsublr{NOAA NCEI}.
  \item الفصل~5 (النماذج الحرارية للغلاف الصخري): عمر قاع المحيط، سيتون وآخرون (\babelsublr{M.~Seton and others}، 2020)، \babelsublr{EarthByte}، \babelsublr{CC~BY~4.0}؛ والتضاريس العالمية \babelsublr{ETOPO1}، \babelsublr{NOAA NCEI}، ملك عام.
  \item الفصل~6 (الحمل في الوشاح والقوى المحرّكة للصفائح): الشواطئ المرفوعة في باليهيلين، مالين هيد، \babelsublr{A.~F.~Borchert}، \babelsublr{CC~BY-SA~4.0} (\babelsublr{Wikimedia Commons})؛ و\babelsublr{PREM}، دزيونسكي وأندرسون (\babelsublr{A.~M.~Dziewonski and D.~L.~Anderson}، 1981).
  \item الفصل~7 (آليات التشوّه ونطاقات القصّ): ميلونيت عيني، \babelsublr{Mikenorton}، \babelsublr{CC~BY-SA~4.0} (\babelsublr{Wikimedia Commons}).
  \item الفصل~9 (التكتونية الامتدادية والانزياحية): تضاريس الأحواض والسلاسل جنوب وادي الموت، كين لوند (\babelsublr{Ken Lund})، \babelsublr{CC~BY-SA~2.0} (\babelsublr{Wikimedia Commons}).
  \item الفصل~14 (جيوكيمياء النظائر المستقرة): المكدَّس القاعي \babelsublr{LR04} لقيم \mbox{\textdir TLT $\delta^{18}$O}، ليسيكي ورايمو (\babelsublr{L.~E.~Lisiecki and M.~E.~Raymo}، 2005).
  \item الفصل~15 (تفاعل الماء والصخر): مصاطب الترافرتين في \babelsublr{Palette Spring}، ينابيع ماموث الحارة، جيم بيكو (\babelsublr{Jim Peaco})، \babelsublr{National Park Service}، ملك عام (\babelsublr{Wikimedia Commons}). ثوابت التوازن من قاعدة بيانات \babelsublr{PHREEQC phreeqc.dat}، باركهيرست وأبيلو (\babelsublr{D.~L.~Parkhurst and C.~A.~J.~Appelo})، \babelsublr{U.S. Geological Survey}، ملك عام.
  \item الفصل~16 (الستراتيغرافيا التتابعية): \babelsublr{Book Cliffs} فوق هيلبر، يوتا، ديفيد ماكي (\babelsublr{David Mackie})، ملك عام (\babelsublr{Wikimedia Commons})؛ والمكدَّس القاعي \babelsublr{LR04} لقيم \mbox{\textdir TLT $\delta^{18}$O}، ليسيكي ورايمو (\babelsublr{L.~E.~Lisiecki and M.~E.~Raymo}، 2005).
  \item الفصل~17 (نماذج السحنات): التطبّق المتقاطع الرابي في تكوين أريسيغ، مايكل ريغل (\babelsublr{Michael C.~Rygel})، \babelsublr{CC~BY-SA~3.0} (\babelsublr{Wikimedia Commons})، مقتطعة.
  \item الفصل~18 (الجيومورفولوجيا الكمية): خانق نياغارا وشلالات حدوة الحصان، كين لوند (\babelsublr{Ken Lund})، \babelsublr{CC~BY-SA~2.0} (\babelsublr{Wikimedia Commons})، مقتطعة.
  \item الفصل~19 (ديناميكا الموائع الجيوفيزيائية): الرياح النطاقية المتوسطة من \babelsublr{NCEP--NCAR Reanalysis~1}، بيانات من \babelsublr{NOAA PSL, Boulder, Colorado, USA}؛ و\babelsublr{World Ocean Atlas 2023}، \babelsublr{NOAA NCEI}.
  \item الفصل~20 (ديناميكا الغلاف الجوي): منخفض جوي قبالة آيسلندا، \babelsublr{NASA/GSFC, MODIS Rapid Response Team}، جاك ديكلواتر (\babelsublr{Jacques Descloitres})، ملك عام (\babelsublr{Wikimedia Commons})؛ والرياح النطاقية المتوسطة من \babelsublr{NCEP--NCAR Reanalysis~1}، بيانات من \babelsublr{NOAA PSL, Boulder, Colorado, USA}.
  \item الفصل~21 (ديناميكا المحيطات): عوّامة \babelsublr{Argo} تُنشر من سفينة أبحاث، \babelsublr{NOAA Global Ocean Monitoring and Observing Program}، ملك عام (\babelsublr{Wikimedia Commons}).
  \item الفصل~22 (نماذج المناخ): حواسيب فائقة في \babelsublr{NASA Center for Climate Simulation}، \babelsublr{Goddard Space Flight Center}، 2010، ملك عام (\babelsublr{Wikimedia Commons})؛ ودرجة حرارة الهواء السطحي والإشعاع عند قمة الغلاف الجوي من \babelsublr{NCEP--NCAR Reanalysis~1}، بيانات من \babelsublr{NOAA PSL, Boulder, Colorado, USA}.
  \item الفصل~24 (علم الجليد): نهر جليدي منفذي في جنوب غرب غرينلاند، \babelsublr{NASA}، جيم يونغل (\babelsublr{Jim Yungel})، \babelsublr{Operation IceBridge}، 2014، ملك عام (\babelsublr{Wikimedia Commons})؛ وارتفاعات سطح غرينلاند من \babelsublr{ETOPO1 Ice Surface}، \babelsublr{NOAA NCEI}؛ وامتدادات الجليد البحري في القطب الشمالي من \babelsublr{NSIDC Sea Ice Index}، الإصدار~4 (\babelsublr{Fetterer and others}، 2025، \babelsublr{doi:10.7265/a98x-0f50}).
  \item الفصل~25 (الاستشعار عن بعد ونظم المعلومات الجغرافية): مخطط التداخل للفترة 1995--2001 لمركز \babelsublr{Three Sisters} البركاني، أوريغون، \babelsublr{U.S. Geological Survey} (\babelsublr{C.~Wicks and others}، 2002)، ملك عام (\babelsublr{Wikimedia Commons})؛ والأطياف المخبرية من \babelsublr{ECOSTRESS Spectral Library}، الإصدار~1.0 (\babelsublr{S.~K.~Meerdink, S.~J.~Hook, D.~A.~Roberts and E.~A.~Abbott}، 2019؛ \babelsublr{A.~M.~Baldridge and others}، 2009)، مستنسخة بإذن من \babelsublr{Jet Propulsion Laboratory, California Institute of Technology, Pasadena, California}.
  \item الفصل~26 (بواطن الكواكب وسطوحها): فوهة بارينجر النيزكية، أريزونا، \babelsublr{USGS/D.~Roddy}، ملك عام (\babelsublr{Wikimedia Commons})؛ وأوليمبوس مونس، فسيفساء المركبة المدارية \babelsublr{Viking}، \babelsublr{NASA/JPL/USGS}، ملك عام (\babelsublr{Wikimedia Commons}). دالتا إنتاج الفوهات وكرونولوجيتها كما في برنامج \babelsublr{Craterstats} لمايكل (\babelsublr{G.~G.~Michael})، عن \babelsublr{G.~Neukum, B.~A.~Ivanov and W.~K.~Hartmann} (2001).
  \item الفصل~27 (الأغلفة الجوية للكواكب وصلاحيتها للحياة): شبكة أودية جافة على المريخ، \babelsublr{Mars Express HRSC}، \babelsublr{ESA/DLR/FU Berlin}، \babelsublr{CC~BY-SA~3.0~IGO} (\babelsublr{Wikimedia Commons})؛ وليجيا ماري على تيتان، فسيفساء رادار \babelsublr{Cassini}، \babelsublr{NASA/JPL/USGS}، ملك عام (\babelsublr{Wikimedia Commons}). ملاءمات النطاق الصالح للحياة لكوبارابو وآخرين (\babelsublr{R.~K.~Kopparapu and others}، 2013)؛ وضغط بخار الميثان من \babelsublr{NIST Chemistry WebBook}.
  \item الفصل~28 (الكواكب الخارجية): \babelsublr{HIP~65426~b} كما صوّره تلسكوب جيمس ويب الفضائي، \babelsublr{NASA/ESA/CSA, A~Carter (UCSC), the ERS~1386 team, and A.~Pagan (STScI)}، \babelsublr{CC~BY~4.0} (\babelsublr{Wikimedia Commons}). الكواكب المؤكَّدة من \babelsublr{NASA Exoplanet Archive}، جدول \babelsublr{Planetary Systems Composite Parameters}، بتاريخ 23 أكتوبر 2026: \babelsublr{this research has made use of the NASA Exoplanet Archive, which is operated by the California Institute of Technology, under contract with the National Aeronautics and Space Administration under the Exoplanet Exploration Program}. منحنيات الكتلة ونصف القطر النموذجية لزينغ وآخرين (\babelsublr{L.~Zeng and others}، 2016 و2019).
  \item الفصل~29 (درب التبانة والمجرّات): بانوراما درب التبانة من مشروع \babelsublr{GigaGalaxy Zoom}، \babelsublr{ESO/S.~Brunier}، \babelsublr{CC~BY~4.0} (\babelsublr{Wikimedia Commons})؛ والمجرّة الحلزونية \babelsublr{M101}، \babelsublr{European Space Agency and NASA} (\babelsublr{Hubble, heic0602; Hubble data K.~D.~Kuntz, F.~Bresolin, J.~Trauger, J.~Mould and Y.-H.~Chu; image processing Davide De~Martin})، \babelsublr{CC~BY~4.0} (\babelsublr{Wikimedia Commons}). السرعات الدائرية في المجرّة من \babelsublr{A.-C.~Eilers, D.~W.~Hogg, H.-W.~Rix and M.~K.~Ness} (2019)؛ ونموذج الكتلة من \babelsublr{E.~Pouliasis, P.~Di~Matteo and M.~Haywood} (2017)؛ والقيفاويات الكلاسيكية من \babelsublr{OGLE} في سحابة ماجلان الكبرى، \babelsublr{I.~Soszy\'nski and others} (2015)، \babelsublr{Acta Astronomica 65, 297}، يوزّعها \babelsublr{CDS, Strasbourg (VizieR J/AcA/65/297)}.
  \item الفصل~30 (من أين تأتي العناصر): سديم السرطان، \babelsublr{NASA, ESA, J.~Hester and A.~Loll (Arizona State University)}، ملك عام (\babelsublr{Wikimedia Commons})؛ وكيلونوفا \babelsublr{GW170817} في \babelsublr{NGC~4993}، \babelsublr{NASA and ESA, acknowledgment A.~J.~Levan, N.~R.~Tanvir, A.~Fruchter and O.~Fox (Hubble, heic1717a)}، \babelsublr{CC~BY~4.0} (\babelsublr{Wikimedia Commons}). الوفرات الشمسية من \babelsublr{M.~Asplund, A.~M.~Amarsi and N.~Grevesse} (2021)؛ ومقاطع الأسر النيوتروني من \babelsublr{KADoNiS v0.3 (I.~Dillmann, R.~Plag, F.~K\"appeler and T.~Rauscher)}؛ وأعمار النصف من \babelsublr{IAEA Live Chart of Nuclides (ENSDF)}؛ ومعدنيات مسح \babelsublr{Geneva--Copenhagen} كما أعاد تحليلها \babelsublr{L.~Casagrande and others} (2011)، يوزّعها \babelsublr{CDS, Strasbourg (VizieR J/A+A/530/A138)}.
\end{itemize}
